%-------------------------------------------------------------------------------
%   ХӨГЖҮҮЛЭЛТЭД АШИГЛАХ АЛГОРИТМ — BikeMap UB
%-------------------------------------------------------------------------------
\section{Хөгжүүлэлтэд ашиглах алгоритм}
\paragraph{} Энэ бүлэг нь системийн хөгжүүлэлтийн гол хэсэг болох \textit{crowd aggregation} (олон хэрэглэгчийн өгөгдлийг нэгтгэх) алгоритм болон \textit{Smart Route} (ухаалаг маршрут тооцоолох) алгоритмыг тайлбарлана. Дэлгэрэнгүй тайлбар, тестийн үр дүнг 2.2.1 -р дэд бүлэгт авч үзнэ.

%-------------------------------------------------------------------------------
\subsection{Crowd Aggregation алгоритм}
\paragraph{}
\textbf{Crowd Aggregation} буюу олон хэрэглэгчийн оруулсан замын нөхцлийн тэмдэглэгээг нэгтгэн тухайн сегментийн \emph{давамгай} (dominant) нөхцлийг тооцоолох алгоритм нь системийн \textit{back-end} хөгжүүлэлтийн гол алгоритм болно.

\textbf{Зарчим:} Сегмент бүрт олон хэрэглэгч 🟢 (\textit{green} — bike lane байгаа), 🟡 (\textit{yellow} — боломжтой), 🔴 (\textit{red} — боломжгүй) гэсэн гурван ангиллын аль нэгийг сонгож тэмдэглэдэг. Систем эдгээр санал (vote) -уудыг нэгтгэж, хамгийн их санал авсан өнгийг тухайн сегментийн \emph{давамгай} нөхцөл болгон газрын зурагт дүрсэлдэг \cite{crowdsource}.

\textbf{Тооцоолох томьёо:}
\begin{equation}
    \text{dominant}(s) = \arg\max_{c \in \{green, yellow, red\}} V_c(s)
\end{equation}
Энд $V_c(s)$ нь $s$ сегментэд $c$ ангилалд өгөгдсөн нийт саналын тоо.

\textbf{Жишээ:} $green = 10$, $yellow = 3$, $red = 6$ бол $\text{dominant} = green$ болж, газрын зурагт ногоон өнгөөр харагдана.

Алгоритмын процесс нь
\begin{enumerate}
    \item Сегментийн \texttt{segment\_hash} -ээр vote-уудыг бүлэглэх
    \item green/yellow/red тус бүрийн санал тоог тоолох
    \item Хамгийн их утгатайг \texttt{dominant} болгож хадгалах
    \item Үр дүнг \texttt{Crowd Aggregation} хүснэгтэд шинэчлэх
\end{enumerate}
гэсэн үндсэн дөрвөн алхамд хуваагддаг.

\begin{figure}[H]
    \centering
    % \includegraphics[scale=1]{Figures/chapter2/crowd-aggregation-architecture.PNG} % auto-commented: missing file
    \caption{Crowd Aggregation алгоритмын архитектур}
    \label{fig:crowdAgg}
\end{figure}

%-------------------------------------------------------------------------------
\subsection{Smart Route — Ухаалаг маршрут тооцоолох алгоритм}
\paragraph{}
\textbf{Smart Route} нь crowd aggregation -ийн үр дүн болон POI (\textit{Point of Interest}) -ийн өгөгдөлд тулгуурлан хэрэглэгчид аюулгүй маршрутыг автоматаар санал болгох алгоритм юм.

\textbf{OSRM /Open Source Routing Machine/} нь OpenStreetMap -н өгөгдөлд тулгуурласан, нээлттэй эхийн маршрут тооцоолох сервер юм \cite{osrm}. Энэхүү сервер нь \texttt{shortest path} (Dijkstra/Contraction Hierarchies) алгоритмыг ашиглан А цэгээс В цэг хүртэлх замыг тооцоолно.

Энгийн маршрут тооцооллоос ялгаатай нь Smart Route нь зам бүрийн \emph{жинг} (weight) дараах хүчин зүйлсээр тохируулна:
\begin{itemize}
    \item Сегментийн давамгай нөхцөл (green=1, yellow=2, red=10)
    \item Тухайн зам дээрх POI-ийн төрөл, тоо
    \item Аюулын POI-ээс 50м-ийн радиуст байгаа эсэх
\end{itemize}

\textbf{Sмart Route -ийн процесс:}
\begin{enumerate}
    \item Хэрэглэгч эхлэл, төгсгөлийн цэгийг сонгох
    \item OSRM серверээс боломжит маршрутуудыг авах
    \item Crowd aggregation -аас сегмент тус бүрийн нөхцлийг авах
    \item Зам бүрд жин (cost) тооцоолох
    \item Хамгийн бага жинтэй замыг сонгож харуулах
\end{enumerate}

\textbf{Маршрутын төрлүүд:}
\begin{itemize}
    \item \textbf{Аюулгүй: }Улаан хэсгийг тойрсон, ногоон/шар сегментүүдийг ашигласан зам.
    \item \textbf{Хурдан: }Хамгийн богино зайтай зам — нөхцөлийг тооцохгүй.
\end{itemize}

\begin{figure}[H]
    \centering
    % \includegraphics[scale=1]{Figures/chapter2/smart-route-architecture.PNG} % auto-commented: missing file
    \caption{Smart Route алгоритмын архитектур}
    \label{fig:smartRoute}
\end{figure}

%-------------------------------------------------------------------------------
\section{Технологийн судалгаа}
\paragraph{} Системийн back-end болон front-end хөгжүүлэлтэд ашиглах технологиуд болон хөгжүүлэлтийг хэрэгжүүлэхэд шаардлагатай хөгжүүлэлтийн орчны тайлбар судалгаа зэргийг агуулсан.

\subsection{Back-End хөгжүүлэлтэд ашиглах технологиуд}

\subsubsection{Django фреймворк}
\par
\textbf{Django} нь 2005 онд анх танилцуулагдсан, BSD License -тэй, Python хэл дээр бичигдсэн өндөр түвшний веб фреймворк юм. \textit{“Batteries-included”} буюу хөгжүүлэлтэд шаардлагатай ихэнх хэрэгслүүдийг агуулсан гэдгээрээ алдартай. Уг фреймворкийг \textit{Django Software Foundation} удирддаг бөгөөд орчин үед Instagram, Pinterest, Mozilla, Disqus гэх мэт томоохон веб платформууд ашиглаж байна \cite{django}.

\par Гол онцлогууд:
\begin{itemize}
    \item \textbf{ORM /Object-Relational Mapping/: }Python класс ашиглан өгөгдлийн санг удирдах боломжтой.
    \item \textbf{Admin Panel: }Автоматаар үүсдэг админ интерфейс нь хэрэглэгчийн удирдлага, POI батлах зэрэг ажилбарыг хялбарчилна.
    \item \textbf{Аюулгүй байдал: }CSRF, XSS, SQL Injection -аас хамгаалах суулгацтай middleware -үүдтэй.
    \item \textbf{Хэмжээ: }Жижиг прототипээс эхлээд том хэмжээний \textit{enterprise} систем хүртэл өргөтгөх боломжтой.
\end{itemize}

\begin{figure}[H]
    \centering
    % \includegraphics[scale=0.4]{Figures/chapter2/django.png} % auto-commented: missing file
    \caption{Django лого}
    \label{fig:django-logo}
\end{figure}

\subsubsection{Django REST Framework (DRF)}
\par
\textbf{Django REST Framework} нь Django -г суурь болгож REST API үүсгэхэд зориулагдсан хүчирхэг сан юм. JSON форматтай өгөгдлийг back-end -ээс front-end рүү дамжуулах, JWT authentication, serializer, pagination зэрэг REST API -д шаардлагатай үндсэн функцуудыг агуулсан \cite{drf}.

\par Гол онцлогууд:
\begin{itemize}
    \item \textbf{Serializers: }Python object-ийг JSON болгож хөрвүүлэх хүчирхэг механизм.
    \item \textbf{ViewSets \& Routers: }CRUD үйлдлийг хэдхэн мөр кодоор бий болгох боломжтой.
    \item \textbf{Browsable API: }Хөгжүүлэгч browser-ээс шууд API -г турших боломжтой.
    \item \textbf{Authentication: }JWT, Token, Session, OAuth зэрэг олон төрлийн нэвтрэлтийг дэмждэг.
\end{itemize}

\subsubsection{Leaflet \& OpenStreetMap}
\par
\textbf{Leaflet} нь 2011 онд гарсан, BSD-2 License -тэй, JavaScript -ээр бичигдсэн интерактив газрын зургийн нээлттэй эхийн сан юм. Хэмжээ нь маш бага (~42KB) бөгөөд гар утсан болон ширээний компьютерт зориулсан responsive газрын зургуудыг үүсгэх боломжийг олгоно \cite{leaflet}.

\par \textbf{OpenStreetMap (OSM)} нь олон нийтийн оролцоонд тулгуурласан газрын зургийн нээлттэй өгөгдлийн сан бөгөөд Google Maps -ийн төлбөртэй API -г орлох гол хувилбар болж байна. Leaflet -тэй хослуулан ашиглах нь BikeMap UB системийн хувьд хэд хэдэн давуу талтай:
\begin{itemize}
    \item \textbf{Үнэгүй: }API key шаардахгүй, хязгааргүй ашиглах боломжтой.
    \item \textbf{Crowd-sourced: }Замын мэдээлэл нь манай төслийн үзэл санаатай ижил, олон нийтийн оролцоогоор бүрддэг.
    \item \textbf{OSRM-тай нийцтэй: }OSM өгөгдлийг шууд маршрут тооцоололд ашиглах боломжтой.
    \item \textbf{Plugin сан: }\texttt{Leaflet.draw}, \texttt{Leaflet.markercluster}, \texttt{Leaflet.heat} гэх мэт олон өргөтгөлтэй.
\end{itemize}

\begin{figure}[H]
    \centering
    % \includegraphics[scale=0.5]{Figures/chapter2/leaflet-osm.png} % auto-commented: missing file
    \caption{Leaflet \& OpenStreetMap лого}
    \label{fig:leaflet-logo}
\end{figure}

\subsubsection{OSRM (Open Source Routing Machine)}
\par
\textbf{OSRM} нь C++ хэл дээр бичигдсэн, өндөр гүйцэтгэлтэй, OpenStreetMap өгөгдлийг ашиглан маршрут тооцооллын үйлчилгээ үзүүлдэг нээлттэй эхийн сервер юм. Дугуй, машин, явган явагчдад зориулсан profile -уудтай бөгөөд BikeMap UB систем нь \texttt{bicycle} profile -ийг ашиглана \cite{osrm}.

\par OSRM нь \textit{Contraction Hierarchies} алгоритмыг ашигладаг бөгөөд хотын хэмжээний маршрутыг хэдхэн миллисекундэд тооцоолж чадна. Нэмж дурдвал custom weight профайл бичих замаар манай Smart Route алгоритмыг хэрэгжүүлэхэд тохиромжтой.

\subsubsection{MySQL өгөгдлийн сан удирдах систем}
\par
MySQL нь бүтэцлэгдсэн өгөгдлийн сангийн менежментийн систем. Анх Майкл Видениус зохион бүтээж байсан бөгөөд анхны хувилбараа 1995 онд гаргасан. 2010 онд тус системийг Oracle компани худалдан авсан, одоо MySQL нь GNU General Public License -ийн нөхцөлийн дагуу үнэгүй, нээлттэй эхийн программ хангамж юм. BikeMap UB системд хэрэглэгчийн мэдээлэл, segment, POI, crowd aggregation хүснэгтүүдийг хадгалахад MySQL -ийг ашиглана.

\par Газарзүйн өгөгдлийг хадгалахдаа MySQL -ийн \texttt{POINT}, \texttt{LINESTRING} зэрэг spatial data type болон \texttt{ST\_Distance\_Sphere()} зэрэг spatial функцуудыг ашиглах боломжтой. Энэ нь POI -ийн ойролцоох сегмент олох, маршрутын радиус тооцоолох зэрэгт хэрэгтэй.

\begin{figure}[H]
    \centering
    % \includegraphics[scale=0.75]{Figures/chapter2/mysql.png} % auto-commented: missing file
    \caption{MySQL лого}
    \label{fig:mysql-logo}
\end{figure}

\subsubsection{PyCharm IDE}
\par
PyCharm нь Python хөгжүүлэгчдэд зориулсан Python -ийн нэгдсэн хөгжүүлэлтийн орчин (IDE) юм. JetBrains компаниар Java болон Python хэл дээр хөгжүүлэгдсэн, 2010 онд анхны хувилбараа гаргасан. Linux, Windows, macOS дээр ажиллах боломжтой бөгөөд нээлттэй эхийн Community болон төлбөртэй Professional хувилбаруудтай. Django integration нь Pro хувилбарт суулгагдсан байдаг бөгөөд URL routing, template debug, Django console зэрэг олон тусгай хэрэгсэлтэй.

\begin{figure}[H]
    \centering
    % \includegraphics[scale=0.9]{Figures/chapter2/pycharm.png} % auto-commented: missing file
    \caption{PyCharm IDE лого}
    \label{fig:pycharmide-logo}
\end{figure}

\par PyCharm IDE -н давуу талууд:
\begin{itemize}
    \item Django, Flask, FastAPI зэрэг Python веб фреймворкуудтай гүн нэгтгэгдсэн.
    \item Pandas, Scikit-learn зэрэг шинжлэх ухааны янз бүрийн номын санд хандах боломжтой.
    \item Database tool window -оор MySQL -тэй шууд холбогдох боломжтой.
    \item Синтаксын алдааг олоход хялбар болгогч кодын алдаа засагчтай.
    \item Автоматаар бөглөх синтакс функц нь цаг хэмнэдэг.
\end{itemize}
PyCharm IDE -н сул талууд:
\begin{itemize}
    \item Сурахад их цаг зарцуулдаг.
    \item Эх кодын хэмжээ томрох тусам RAM -ийн хэрэглээ нэмэгддэг.
    \item Олон нийтийн Community хувилбар нь Django -ийн зарим тусгай функцийг дэмждэггүй.
\end{itemize}

%-------------------------------------------------------------------------------
\subsection{Front-End хөгжүүлэлтэнд ашиглах технологиуд}

\subsubsection{React фреймворк}
React (\textit{ReactJS, React.js} гэж бас хэлэгддэг) нь веб аппликейшны \textit{User Interface} (хэрэглэгчийн интерфейс) -г хөгжүүлэхэд зориулагдсан JavaScript фреймворк юм. 2011 онд Мета (Facebook) компанийн программ хангамжийн инженер Жордан Уолке JavaScript хөгжүүлэлтийн нарийн төвөгтэй байдлыг арилгахын тулд \textit{components} (бүрэлдэхүүн хэсгүүд) -д суурилсан, дахин ашиглах боломжтой React фреймворкыг хөгжүүлж, 2013 онд анхны хувилбараа гаргаж байсан \cite{react}.

\par BikeMap UB системийн хувьд интерактив газрын зураг, бодит цагийн POI vote, маршрут визуалчлал зэрэг динамик UI элементүүдийг гаргахад React тохиромжтой. Мөн \texttt{react-leaflet} сан нь Leaflet -ийг React component хэлбэрээр ашиглах боломжийг олгодог.

\par ReactJS -н давуу талууд:
\begin{itemize}
    \item Суралцах болон хэрэглэхэд хялбар. Ойлгоход хялбар документ болон зааварчилгаатай. Хамгийн олон практик хичээлтэй фреймворк.
    \item JSX (JavaScript Extension) -г ашиглаж динамик веб хөгжүүлэхэд хялбар болгосон.
    \item Virtual DOM ашигласнаар гүйцэтгэлийг нэмэгдүүлсэн.
    \item Дахин ашиглах боломжтой component -г бий болгодог.
    \item \texttt{react-leaflet}, \texttt{axios}, \texttt{react-router} зэрэг өргөн экосистемтэй.
\end{itemize}

ReactJS -н сул талууд:
\begin{itemize}
    \item Байнгын шинэчлэл. Энэ нь хөгжүүлэгчдийн хувьд шинэчлэлтийг байнга судалж, хөгжүүлэлтэд нэвтрүүлэх шаардлагатай байдаг учраас хөгжүүлэлтийг удаашруулах хандлагатай байдаг.
    \item Системийн хувьд зөвхөн хэрэглэгчийн интерфейсийн хэсгийг хамардаг учраас системийг бүхэлд нь хөгжүүлэхийн тулд бусад технологийг ч бас нэвтрүүлэх шаардлагатай гардаг.
\end{itemize}

\begin{figure}[H]
    \centering
    % \includegraphics[scale=0.055]{Figures/chapter2/react.png} % auto-commented: missing file
    \caption{React лого}
    \label{fig:react-logo}
\end{figure}

\subsubsection{Bootstrap 5}
\par
\textbf{Bootstrap} нь Twitter компаний инженерүүд Mark Otto, Jacob Thornton нар 2011 онд боловсруулсан, MIT License -тэй, нээлттэй эхийн front-end фреймворк юм. Responsive (\textit{mobile-first}) дизайны зарчмыг үндэс болгож хөгжүүлэгдсэн \cite{bootstrap}. BikeMap UB системийн хувьд гар утас болон ширээний компьютерт ижилхэн харагдах хэрэглэгчийн интерфейс үүсгэхэд Bootstrap 5 ашиглагдана.

\par Bootstrap -ийн онцлогууд:
\begin{itemize}
    \item \textbf{Grid system: }12-баганат уян хатан grid system нь responsive layout үүсгэхэд тохиромжтой.
    \item \textbf{Бэлэн компонент: }Button, Modal, Form, Card, Navbar гэх мэт олон төрлийн UI компонент.
    \item \textbf{Utility class: }Margin, padding, color гэх мэт өөрчлөлтийг class -аар хийх боломжтой.
    \item \textbf{Икон сан: }\texttt{Bootstrap Icons} нь 1800+ үнэгүй икон агуулсан.
    \item \textbf{Документ: }Дэлхийн хамгийн их жишээтэй фреймворкуудын нэг.
\end{itemize}

\begin{figure}[H]
    \centering
    % \includegraphics[scale=0.5]{Figures/chapter2/bootstrap.png} % auto-commented: missing file
    \caption{Bootstrap 5 лого}
    \label{fig:bootstrap-logo}
\end{figure}

\subsubsection{Visual Studio Code}
\par
\textbf{Visual Studio Code} (VS Code гэж нэрлэдэг) нь Microsoft компани Electron Framework -тэй хамтран хийсэн, Windows, Linux, macOS -д зориулсан кодын засварлагч юм. 2015 оны 11-р сарын 18 -нд VS Code -ийн эх код MIT лицензийн дагуу гарч, GitHub дээр тавигдсан. Онцлогуудад нь дибаг хийх, синтакс онцлох, ухаалаг код гүйцээлт, кодын дахин боловсруулалт, суулгагдсан Git зэрэг орно. Хэрэглэгчид загвар, гарын товчлол, тохиргоог өөрчлөх, нэмэлт функц бүхий өргөтгөлүүдийг суулгах боломжтой.

\par BikeMap UB системийн front-end хөгжүүлэлтэд React snippets, ESLint, Prettier, GitLens зэрэг өргөтгөлүүд ашиглагдана. Мөн \texttt{Live Server} өргөтгөлөөр локал хөгжүүлэлтийн серверийг гаргах боломжтой.

\begin{figure}[H]
    \centering
    % \includegraphics[scale=0.8]{Figures/chapter2/vscode.png} % auto-commented: missing file
    \caption{Visual Studio Code лого}
    \label{fig:vscode-logo}
\end{figure}

%-------------------------------------------------------------------------------
\subsection{Кодын менежмент}
\par BikeMap UB системийн эх кодыг \textbf{GitHub} платформ дээр хадгалж, \textit{Git} систем ашиглан хувилбар удирдлагыг гүйцэтгэнэ. Хөгжүүлэлтийн ажилтай уялдуулан \textit{branching strategy} -ийг дараах байдлаар ашиглана:
\begin{itemize}
    \item \texttt{main} — production хувилбарын код
    \item \texttt{develop} — идэвхтэй хөгжүүлэлтийн код
    \item \texttt{feature/*} — шинэ функц хөгжүүлэх branch (жш: \texttt{feature/poi-vote})
    \item \texttt{fix/*} — алдаа засах branch
\end{itemize}

\par GitHub -ийн \textit{Pull Request} механизм нь кодын чанарыг хянах, code review хийх боломжийг олгоно. Мөн \textbf{GitHub Actions} ашиглан CI/CD pipeline -ийг автоматжуулна — push хийх бүрд unit test ажиллах, lint шалгах, deployment хийх зэрэг ажилбар автомат гүйцэтгэгдэнэ.

%-------------------------------------------------------------------------------
\section{Бүлгийн дүгнэлт}
%-------------------------------------------------------------------------------
\paragraph{} Тус бүлэгт BikeMap UB системийн гол алгоритм болох \textbf{Crowd Aggregation} болон \textbf{Smart Route} -ийн ажиллах зарчмыг тайлбарласан. Crowd aggregation нь олон хэрэглэгчийн оруулсан саналыг нэгтгэн сегмент бүрийн давамгай нөхцлийг тооцоолох энгийн ч найдвартай арга юм. Жишээ нь нэг сегментэд $green=10$, $yellow=3$, $red=6$ санал орвол алгоритм нь $\arg\max$ үйлдлээр \emph{green} -ийг давамгай нөхцөл болгож тогтооно. Smart Route алгоритм нь OSRM серверээс авсан маршрут дээр crowd aggregation -ийн үр дүн болон POI -ийн өгөгдлийг хослуулан зам бүрд жин (cost) тооцоолж, аюулгүй маршрутыг автоматаар санал болгоно.

\paragraph{} Тус бүлэгт системийг хөгжүүлэхэд шаардлагатай технологиудыг тайлбарласан. Back-end хөгжүүлэлтэд \textbf{Django} веб фреймворк, REST API үүсгэхэд \textbf{Django REST Framework}, өгөгдлийн санд \textbf{MySQL} (spatial data type-тэй), маршрут тооцооллын серверт \textbf{OSRM} -ийг ашиглана. Газрын зураг ба интерактив элементүүдэд \textbf{Leaflet} болон \textbf{OpenStreetMap} -ийг хослуулан ашиглах нь үнэ төлбөргүй, манай системийн crowd-sourcing үзэл санаатай нийцэх давуу талтай. Front-end хөгжүүлэлтэнд \textbf{React} фреймворк, responsive дизайнд \textbf{Bootstrap 5}, кодын засварлагчид \textbf{PyCharm} (back-end) болон \textbf{Visual Studio Code} (front-end) -ийг ашиглана. Эх кодын менежментийг \textbf{GitHub} дээр Git ашиглан хийх бөгөөд GitHub Actions -оор CI/CD pipeline-ийг автоматжуулна.
